\documentclass[11pt]{book}

\usepackage{ppbook}

\begin{document}

\setcounter{chapter}{5}
\chapter{Automated Patch Synthesis}

APR = automated program repair.

Basic flow:

\[
\text{fault localisation}
\rightarrow
\text{patch synthesis}
\rightarrow
\text{verification}
\]

First find where the bug probably is,
then generate some possible fixes,
then check whether the fix actually works.

For a buggy program $P$, we want some patch $y$
such that the patched program satisfies the property.

\[
\exists y\;.\;\forall x\;.\;
\phi(P\oplus y,x)
\]

So patch synthesis is basically finding a $y$
which works for all inputs.

\section{CEGIS}

Patch problem fits naturally into
counterexample-guided inductive synthesis.

Don't try all possible inputs initially.

Start with a finite set:

\[
X=\{x_1,x_2,\ldots,x_n\}
\]

Find a patch satisfying these examples:

\[
\exists y\;.\;
\bigwedge_i \phi(P\oplus y,x_i)
\]

Then the verifier checks whether there is some input
for which the patch still fails:

\[
\exists x\;.\;
\neg\phi(P\oplus y,x)
\]

If such an $x$ is found,

\[
X \leftarrow X\cup\{x\}
\]

and synthesize again.

So:

\[
\text{examples}
\rightarrow
\text{patch}
\rightarrow
\text{counterexample}
\rightarrow
\text{more examples}
\]

\section{How patches are searched}

One way = grammar.

Define possible expressions / replacements,
then search through them.

Example grammar can contain

\[
age>c,\qquad age\geq c,\qquad age<c
\]

and use one of these to replace the hole.

Another way = mutation.

Operations like

\[
delete,\quad insert,\quad replace
\]

are applied to the existing program.

Constraint-based synthesis is another option.

There is a hole in the program:

\[
HOLE(x)
\]

and synthesis searches for an expression that satisfies
the required constraint.

\section{Patch size}

Usually prefer a small patch.

Give a cost to a patch:

\[
C=\alpha\cdot LinesChanged+
\beta\cdot Complexity
\]

Why?

If there is no such preference,
the repair system could just rewrite a large part
of the program.

\begin{ppnote}
Minimal repair assumption:

the intended program is probably close to the
buggy program.
\end{ppnote}

For one line there are roughly $N$ alternatives,
but for $k$ lines the search space becomes much bigger.

Stateful systems are even harder because correctness
can depend on

\[
\text{state + ordering + external calls + exceptions}
\]

\section{Finding the bug}

A failed assertion does not necessarily mean
the exact line near the assertion is wrong.

The problem could come from:

\begin{itemize}
    \item caller
    \item callee
    \item sequence of statements
\end{itemize}

So we need fault localisation.

Some techniques:

\begin{itemize}
    \item spectrum based fault localisation
    \item dynamic slicing
    \item symbolic execution
    \item mutation analysis
    \item LLM reasoning
\end{itemize}

\section{LLM-generated code}

LLM-generated programs can also contain bugs,
so APR can be applied to them.

Some bugs are small changes:

\[
\text{operator change}
\]

\[
\text{variable replacement}
\]

etc.

Other bugs need larger patches.

A significant type of bug is algorithmic misalignment,
where the generated code does not implement the
intended algorithm correctly.

Traditional APR methods can use statistical fault
localisation to identify suspicious lines.

But just knowing suspicious lines is not always enough.

Program dependencies matter too.

A suspicious statement may depend on another statement
somewhere else.

\section{LLM as repairer}

Instead of APR generating the patch itself,
the LLM can be asked to repair the program.

Different levels of giving information:

\[
\text{whole program}
\]

or

\[
\text{program + suspicious lines}
\]

or even

\[
\text{program + particular statement to fix}
\]

Giving location information can guide the model,
but the generated patch still needs verification.

\section{Combining APR and LLMs}

TBar / Recoder type APR methods can be combined
with Codex-generated solutions.

Instead of relying on one generated patch,
multiple candidate solutions can be produced.

Then patch ingredients from the candidates can be
used together with APR.

Rough idea:

\[
\text{APR patches}
+
\text{LLM patches}
\rightarrow
\text{more candidate fixes}
\]

and then verify them.

The useful direction is not just

\[
\text{LLM}\rightarrow\text{code}
\]

but

\[
\text{LLM}
\leftrightarrow
\text{program repair}
\]

APR can give structure / fixing operators,
while LLMs can provide different possible repairs.

Domain knowledge can also help in learning better
fix patterns and make those patterns more general.

\begin{ppkey}
The overall repair loop is

\[
\boxed{
\text{locate bug}
\rightarrow
\text{generate patches}
\rightarrow
\text{verify}
\rightarrow
\text{counterexample if needed}
}
\]
\end{ppkey}

Main difficulty = huge patch search space + finding
the actual fault location + making sure the patch
really fixes the intended behaviour.

\end{document}